\scanpage{5}
\mspage{4}
EGA IV 18.9.\qquad (\uline{Rabiots})

L'assertion de connexité faite dans 18.9.1. se généralise de la façon suivante:

\uline{Théorème} 18.9.7. Soit $A\quad B$ un homomorphisme local d'anneaux locaux noethériens, et supposons a) \inserted{$f:$} $\operatorname{Spec}(B)=X\quad\operatorname{Spec}(A)=Y$ est plat à fibres séparables, et b) $A$ est géométriquement unibranche, ou $A$ est unibranche et $k(B)$ une extension primaire de $k(A)$. Alors la fibres générique de $X$ sur $Y$ est connexe, et pour toute partie fermée $S$ de $Y$ distincte de $Y$, \deleted{\unsure{et} \illegible{1}} et tout fermé $T$ de $X$ contenu dans $f^{-1}(S)$, $X-T$ est connexe.

Signalons tout \unsure{des|de} suite les corollaires suivants:

\uline{Corollaire} 18.9.8. Supposons qu'on ait a) ci-dessus, et b') $A$ est strictement local, ou $A$ est local hensélien et $k(B)$ est une extension primaire de $k(A)$. Alors \uline{toutes} les fibres de $X$ sur $Y$ sont \uline{géométriquement connexes}.

En effet, soit $A\quad A'$ un homomorphisme local, avec $A'$ fini sur $A$, et soit $B'=B\otimes_A A'$, alors on constate aussitôt que les hypothèses faites sont encore vérifiées pour $A'\quad B'$. Ceci nous ramène, comme d'habitude, au cas où $A$ est intègre, et à vérifier que dans ce cas, la fibre générique de $f:X\quad Y$ est connexe. Mais nous sommes alors sous les conditions du théorème qui précède.

\uline{Corollaire} 18.9.9. Soit $A$ un anneau local noethérien hensélien, \deleted{\unsure{int}} universellement caténaire (i.e. réf\quad à fibres formelles séparables). Alors les fibres formelles de $A$, i.e. les fibres de $\operatorname{Spec}(\hat{\hat A})\quad\operatorname{Spec}(A)$, sont \inserted{géom.} connexes.

C'est en effet un cas particulier de 18.9.8. (l'extension résiduelle $k(\hat A)$ sur $k(A)$ est ici triviale, à fortiori primaire).
